\documentclass[12pt,a4paper]{article}
\usepackage[utf8]{inputenc}
\usepackage[T1]{fontenc}
\usepackage[bahasa]{babel}
\usepackage{mathptmx} % Times New Roman font
\usepackage{microtype} % Untuk optimasi typography dan spacing
\usepackage[none]{hyphenat} % Menonaktifkan hyphenation
\usepackage{graphicx}
\usepackage{amsmath}
\usepackage{amssymb}
\usepackage{setspace}
\usepackage{geometry}
\usepackage{tikz}
\usepackage{float}
\usepackage{enumitem}
\usepackage{longtable}
\usepackage{tabularx}
\usepackage{array}
\usepackage{hyperref}
\hypersetup{
    pdfencoding=auto,
    pdftitle={Daftar Lambang dan Singkatan},
    pdfauthor={},
    pdfsubject={Tesis},
    colorlinks=true,
    linkcolor=black,
    citecolor=black,
    urlcolor=blue,
    pdfstartview=FitH,
    bookmarksopen=true,
    bookmarksnumbered=true
}
\usepackage{tocloft}
\usepackage{caption}
\usepackage{textcomp} % For special symbols


% Aktifkan hyphenation Bahasa Indonesia dan optimasi spasi agar tidak keluar margin
\lefthyphenmin=2
\righthyphenmin=2
\pretolerance=1000
\tolerance=2000
\emergencystretch=3em
\hbadness=10000
\vbadness=10000
\hyphenpenalty=500
\exhyphenpenalty=500
\sloppy
% Aktifkan microtype untuk Bahasa Indonesia
\microtypesetup{activate=true,protrusion=true,expansion=true,tracking=true,kerning=true}
\microtypecontext{spacing=nonfrench}

% Pengaturan paragraf: tanpa indentasi dan jarak minimal antar paragraf
\setlength{\parindent}{0pt} % Menghilangkan indentasi awal paragraf
\setlength{\parskip}{0pt} % Tidak ada jarak antar paragraf

% Pengaturan jarak untuk subsection
\usepackage{titlesec}
\titleformat{\subsection}{\normalsize\bfseries}{\thesubsection}{1em}{}
\titlespacing*{\subsection}{0pt}{0pt}{0pt} % Tidak ada jarak sebelum atau setelah judul

% Pengaturan margin: kiri 4cm, lainnya 3cm
\geometry{
    a4paper,
    left=4cm,
    right=3cm,
    top=3cm,
    bottom=3cm
}

% Pengaturan spasi
\singlespacing

% Pengaturan penomoran halaman dengan angka Romawi kecil
\pagenumbering{roman}

% Pengaturan penomoran
\setcounter{secnumdepth}{3}
\setcounter{tocdepth}{3}

% Pengaturan heading
\renewcommand{\thesection}{\Roman{section}}
\renewcommand{\thesubsection}{\thesection.\arabic{subsection}}
\renewcommand{\thesubsubsection}{\thesubsection.\arabic{subsubsection}}

\begin{document}

% ============================================================
% DAFTAR LAMBANG DAN SINGKATAN
% ============================================================

\vspace{2cm}
\begin{center}
{\fontsize{14}{16.8}\selectfont\textbf{DAFTAR SINGKATAN DAN LAMBANG}}\\[1em]
\end{center}
\label{sec:lambang-singkatan}
\addcontentsline{toc}{section}{DAFTAR LAMBANG DAN SINGKATAN}
\setcounter{section}{0}
\setcounter{subsection}{0}
\vspace{2em}

% ============================================================
% SECTION 1: DAFTAR SINGKATAN
% ============================================================

% \subsection{Daftar Singkatan}
\label{subsec:daftar-singkatan}
\vspace{0.8em}

\begin{longtable}{@{} l >{\raggedright\arraybackslash}p{7.5cm} c @{}}
\multicolumn{1}{l}{\textbf{SINGKATAN}} & \multicolumn{1}{l}{\textbf{Nama}} & \multicolumn{1}{l}{\parbox[c]{3.5cm}{\raggedright \textbf{Pemakaian pertama kali pada halaman}}} \\

% ============================================================
% DATA SINGKATAN (ALFABETIS A-Z)
% ============================================================

AI & \textit{Artificial Intelligence} & 2 \\
ANRI & Arsip Nasional Republik Indonesia & 1 \\
API & \textit{Application Programming Interface} & 85 \\
BCE & \textit{Binary Cross-Entropy} & 30 \\
BiGRU & \textit{Bidirectional Gated Recurrent Unit} & 25 \\
BiLSTM & \textit{Bidirectional Long Short-Term Memory} & 25 \\
CBAM & \textit{Convolutional Block Attention Module} & 35 \\
CER & \textit{Character Error Rate} & 5 \\
cGAN & \textit{Conditional Generative Adversarial Network} & 20 \\
CLI & \textit{Command-Line Interface} & 90 \\
CNN & \textit{Convolutional Neural Network} & 15 \\
CRNN & \textit{Convolutional Recurrent Neural Network} & 20 \\
CTC & \textit{Connectionist Temporal Classification} & 25 \\
CUDA & \textit{Compute Unified Device Architecture} & 90 \\
cuDNN & \textit{CUDA Deep Neural Network library} & 90 \\
CWV & \textit{Core Web Vitals} & 95 \\
DE-GAN & \textit{Document Enhancement GAN} & 18 \\
DIBCO & \textit{Document Image Binarization Contest} & 22 \\
DocEnTr & \textit{Document Enhancement Transformer} & 19 \\
DRD & \textit{Distance Reciprocal Distortion} & 22 \\
DSRM & \textit{Design Science Research Methodology} & 50 \\
ERB & \textit{Enhance to Read Better} & 18 \\
FCN & \textit{Fully Convolutional Network} & 16 \\
FM & \textit{F-measure} & 22 \\
Fps & \textit{pseudo-F-measure} & 22 \\
FR & \textit{Functional Requirement} & 60 \\
GAN & \textit{Generative Adversarial Network} & 2 \\
GPU & \textit{Graphics Processing Unit} & 85 \\
GT & \textit{Ground Truth} (Kebenaran Dasar) & 55 \\
H-DIBCO & \textit{Handwritten Document Image Binarization Contest} & 22 \\
HMM & \textit{Hidden Markov Model} & 25 \\
HTR & \textit{Handwritten Text Recognition} & 3 \\
IAM & \textit{IAM Handwriting Database} & 27 \\
ICDAR & \textit{International Conference on Document Analysis and Recognition} & 22 \\
IIIT5K & \textit{IIIT 5K-word dataset} & 27 \\
KHATT & \textit{Arabic Handwritten Text Database} & 27 \\
LSTM & \textit{Long Short-Term Memory} & 25 \\
MAE & \textit{Mean Absolute Error} & 30 \\
ML & \textit{Machine Learning}  & 2 \\
MLflow & \textit{Machine Learning Flow} & 90 \\
MSE & \textit{Mean Squared Error}  & 30 \\
MSFP & \textit{Multi-Scale Feature Pyramid} & 35 \\
NFR & \textit{Non-Functional Requirement}  & 65 \\
OCR & \textit{Optical Character Recognition} & 3 \\
PALM & \textit{Presentation Attack Liveness Detection in Mobile scenarios} & 27 \\
Pix2Pix & \textit{Image-to-Image Translation with Conditional Adversarial Networks} & 19 \\
PSNR & \textit{Peak Signal-to-Noise Ratio} & 5 \\
RDB & \textit{Residual Dense Block} & 35 \\
RecFeat & \textit{Recognition Feature Loss} & 120 \\
RMSProp & \textit{Root Mean Square Propagation} & 85 \\
RNN & \textit{Recurrent Neural Network}  & 25 \\
SDM & Sumber Daya Manusia & 60 \\
SGD & \textit{Stochastic Gradient Descent} & 85 \\
SNR & \textit{Signal-to-Noise Ratio}  & 30 \\
SOP & \textit{Standard Operating Procedure}  & 65 \\
SOTA & \textit{State-of-the-Art} & 5 \\
SSIM & \textit{Structural Similarity Index Measure} & 5 \\
UNESCO & \textit{United Nations Educational, Scientific and Cultural Organization} & 1 \\
U-Net & \textit{U-shaped Convolutional Network} & 16 \\
VGG & \textit{Visual Geometry Group} & 17 \\
ViT & \textit{Vision Transformer} & 19 \\
VOC & \textit{Visual Object Classes} & 27 \\
WER & \textit{Word Error Rate} & 28 \\

\end{longtable}

\newpage

% ============================================================
% SECTION 2: DAFTAR LAMBANG DAN SIMBOL MATEMATIS
% ============================================================

% \subsection{Daftar Lambang dan Simbol Matematis}
\label{subsec:daftar-lambang}
\vspace{0.8em}

\begin{longtable}{@{} l >{\raggedright\arraybackslash}p{6.5cm} c @{}}
\multicolumn{1}{l}{\textbf{LAMBANG}} & \multicolumn{1}{l}{\textbf{Nama}} & \multicolumn{1}{l}{\parbox[c]{3.5cm}{\raggedright \textbf{Pemakaian pertama kali pada halaman}}} \\

% ============================================================
% DATA LAMBANG DAN SIMBOL (KATEGORIS)
% ============================================================

% \multicolumn{2}{l}{\textbf{Huruf Yunani (Parameter Model)}} \\

$\alpha$ & Koefisien bobot untuk \textit{loss function}, tingkat signifikansi statistik & 75 \\
$\beta$ & Koefisien bobot untuk \textit{Binary Cross-Entropy loss} & 75 \\
$\lambda$ & Koefisien bobot untuk CTC \textit{loss}, parameter regularisasi & 75 \\
$\theta$ & Parameter model \textit{neural network} & 75 \\
$\sigma$ & Deviasi standar & 75 \\
$\mu$ & Rata-rata (\textit{mean}) & 75 \\
$\epsilon$ & \textit{Token blank}/kosong dalam CTC, konstanta kecil untuk stabilitas numerik & 75 \\
$\Delta$ & Delta (selisih/perubahan nilai) & 75 \\
$\pi$ & \textit{Alignment path} dalam CTC & 75 \\
$\pi_t$ & \textit{Token} pada \textit{timestep} $t$ dalam \textit{alignment path} & 76 \\[1ex]

% \multicolumn{2}{l}{\textbf{Operator Matematis}} \\

$\nabla$ atau $\partial$ & Operator gradien atau turunan parsial & 75 \\
$\sum$ & Operator penjumlahan & 75 \\
$\prod$ & Operator perkalian & 75 \\
$\int$ & Operator integral & 75 \\
argmax & Argumen yang memaksimalkan fungsi & 76 \\
$\sim$ & Terdistribusi menurut & 75 \\
$\to$ & Menuju atau konvergen ke & 75 \\[1ex]

% \multicolumn{2}{l}{\textbf{Fungsi \textit{Loss} dan Kerugian}} \\

$\mathcal{L}$ & Fungsi \textit{loss} (kerugian) & 75 \\
$\mathcal{L}_{total}$ & Total \textit{loss function} & 75 \\
$\mathcal{L}_{adversarial}$ atau $\mathcal{L}_{adv}$ & \textit{Adversarial loss} & 75 \\
$\mathcal{L}_{reconstruction}$ & \textit{Reconstruction loss} & 75 \\
$\mathcal{L}_{L1}$ & L1 \textit{loss} (\textit{Mean Absolute Error}) & 75 \\
$\mathcal{L}_{L2}$ & L2 \textit{loss} (\textit{Mean Squared Error}) & 75 \\
$\mathcal{L}_{CTC}$ & CTC (\textit{Connectionist Temporal Classification}) \textit{loss} & 75 \\
$\mathcal{L}_{perceptual}$ & \textit{Perceptual loss} & 75 \\
$\mathcal{L}_{pixel}$ & \textit{Pixel-wise loss} & 75 \\
$\mathcal{L}_{BCE}$ & \textit{Binary Cross-Entropy loss} & 75 \\
$\mathcal{L}_{GAN}$ & GAN \textit{loss} & 75 \\
$\mathcal{L}_{cycle}$ & \textit{Cycle consistency loss} & 75 \\[1ex]

% \multicolumn{2}{l}{\textbf{Variabel Model dan Komponen GAN}} \\

$G$ & \textit{Generator} dalam GAN & 75 \\
$D$ & \textit{Discriminator} dalam GAN & 75 \\
$G(z)$ & \textit{Output} dari \textit{generator} dengan \textit{input} $z$ & 75 \\
$D(x)$ & \textit{Output} dari \textit{discriminator} dengan \textit{input} $x$ & 75 \\
$G_{A \to B}$ & \textit{Generator} dari domain $A$ ke $B$ & 75 \\
$G_{B \to A}$ & \textit{Generator} dari domain $B$ ke $A$ & 75 \\
$V(D, G)$ & Fungsi nilai (\textit{value function}) dalam GAN & 75 \\
$x$ & Sampel data asli, \textit{input} citra & 75 \\
$y$ & Label atau \textit{ground truth}, kondisi dalam \textit{conditional GAN} & 75 \\
$z$ & \textit{Noise vector} atau representasi laten \\[1ex]

% \multicolumn{2}{l}{\textbf{Distribusi Probabilitas}} \\

$\mathbb{E}$ & Ekspektasi atau nilai harapan & 75 \\
$p_{data}$ & Distribusi data asli & 75 \\
$p_z$ & Distribusi \textit{prior} (\textit{noise}) & 75 \\
$p(y|x)$ & Probabilitas bersyarat \textit{output} $y$ \textit{given input} $x$ & 75 \\
$p_t(\pi_t|x)$ & Probabilitas \textit{token} pada \textit{timestep} $t$ \\[1ex]

% \multicolumn{2}{l}{\textbf{Fungsi dan Transformasi CTC}} \\

$\mathcal{B}$ & Fungsi penciutan (\textit{collapse function}) dalam CTC & 75 \\
$\mathcal{B}^{-1}$ & Fungsi invers penciutan dalam CTC & 75 \\
$T$ & Jumlah \textit{timestep} atau panjang \textit{sequence} \\[1ex]

% \multicolumn{2}{l}{\textbf{Metrik Jarak dan Norma}} \\

$|x - G(y)|_1$ & L1 \textit{distance} antara \textit{ground truth} dan \textit{generated image} & 75 \\
$|x - G(y)|_2$ & L2 \textit{distance} antara \textit{ground truth} dan \textit{generated image} & 75 \\
$N \times N$ & Ukuran \textit{patch} dalam PatchGAN \\[1ex]

% \multicolumn{2}{l}{\textbf{Statistik dan Analisis}} \\

$d$ & Cohen's $d$ (\textit{effect size}) & 75 \\
$n$ & Jumlah sampel & 75 \\
$p$ & \textit{p-value} (nilai probabilitas statistik) & 75 \\
$r$ & Koefisien korelasi & 75 \\
$R^2$ & Koefisien determinasi \\[1ex]

% \multicolumn{2}{l}{\textbf{Notasi Kompleksitas Komputasi}} \\

$O(n)$ & Notasi Big-O untuk kompleksitas linear & 75 \\
$O(n^2)$ & Notasi Big-O untuk kompleksitas kuadratik & 75 \\

\end{longtable}

\end{document}
